% Proposed replacement for the Discussion in the v2 manuscript.
% Recheck numerical claims against the completed v2 campaign before integration.
\section{Discussion}
\label{sec:discussion}

The simulations support a discontinuous active--absorbing transition in a two-dimensional
binary model with exactly two symmetry-related absorbing configurations. At $p=1/2$, the
front reversal occurs while the active phase retains a finite interface density.
The active phase retreats as a coherent region, and bulk runs on its metastable side
decay through the growth of absorbing domains. Large droplets of the competing phases
exchange growth and shrinkage across the same narrow parameter range. The orientation
tests constrain a broad interval of generic coexistence~\cite{wang2015,liu2018}.
Together these measurements identify phase competition through independently measured
bulk and interfacial observables.

This result extends the two-dimensional picture of GV criticality and
Ising--DP splitting~\cite{dornic2001,droz2003,alhammal2005}. Earlier reports of discontinuous
behavior in two-state models~\cite{rodrigues2015} make the distinction between activity and
magnetization essential. The contribution here is the combined evidence for coexistence
in a model whose absorbing configurations are explicitly characterized.
The front phenomenology has close analogues in defensive-alliance
models~\cite{szabo2004,guisoni2013}; realizing it with binary sites and only two absorbing
configurations isolates the transition from internal dynamics within an absorbing phase.

A physical interpretation is competition between active and absorbing regions,
with front propagation opposed by curvature. On the side where an absorbing domain
invades, a sufficiently large fluctuation can initiate decay of the active phase.
On the other side, small absorbing domains contract.
A predictive theory would need to derive the front mobility, curvature response,
and nucleation rate from the microscopic rule. A local deterministic potential and
its noise amplitude alone do not determine those spatial quantities.

The continuation toward the voter point introduces a separate scale problem.
A small density jump, slow continuous vanishing, and a long crossover can be difficult
to distinguish over finite times. Nonlinear GV models can also have transient
magnetization drift and seed dynamics that differ from the linear voter
model~\cite{castellano2012,azizi2020}. Effective exponents therefore support the
interpretation together with independent boundary estimates and converged bulk and
front measurements. The asymptotic endpoint remains an open question.

The coalescing-walk closure organizes the microscopic perturbations and accounts
approximately for the scale $q^*\sim\alpha(1-p)$. Higher-lineage contributions supply
terms capable of bistability within the closure, and their decreasing weight suggests
a narrowing bistable region. Establishing its asymptotic width requires a controlled
matching scale and the fluctuations beyond the local reaction term. Connecting this
construction to rigorous two-dimensional perturbation limits~\cite{cox2010,cox2025}
remains open.

The five-cell comparison makes the limitation concrete. The closure places the invasion
boundary close to simulation, but no single matching time describes that boundary and
all of its predicted stability limits. In particular, simulations do not resolve its
predicted metastable active branch. Small-seed suppression and a finite-time droplet
threshold characterize invasion dynamics; they do not alone settle the limiting order
of the transition. Understanding why the Day \& Night construction supports the
stronger coexistence signatures is a target for a theory of nonlinear voter rules.

A systematic weak-noise study of the spatial Langevin equation would test which
continuum parameters reproduce the observed coexistence. A controlled approach to the
voter endpoint could then establish whether the discontinuity survives arbitrarily
close to it or terminates before it.
